SGA 4 XVI 3.7), and, conjecturally, when @@M0@@ and @@M1@@ are both regular (by the absolute cohomological purity theorem, proved at least for excellent @@M2@@ of characteristic zero (SGA 4 XIX 3.2)). When @@M3@@ and @@M4@@ are smooth over the field of complex numbers, one can also use the integral Chern class on the locally compact space @@M5@@ of complex points of @@M6@@. Formula (3.1) is established in this case by \uline{Atiyah} and \uline{Hirzebruch} [1], who work more generally with an immersion of complex analytic varieties (using the ``naive'' @@M7@@); it follows that the @@M8@@-adic formula (3.1) holds whenever @@M9@@ and @@M10@@ are smooth over an \uline{algebraically closed field of characteristic zero}, using the ``Lefschetz principle'' and the general results of SGA 4. This therefore makes the validity of the cohomological formula (3.1) extremely plausible; but observe that it has not been proved at present even when @@M11@@ and @@M12@@ are smooth over a field @@M13@@ of characteristic zero, when @@M14@@ is not algebraically closed.

\sgaStableTarget{sga6.e14.para.3-4}\textbf{3.4.} One can also work with the cohomological theory of Chern classes in the De Rham or Hodge style [13]. The question of the validity of (3.1) in this setting arises only when @@M15@@ is not of characteristic zero, since in any case (no.~6) the formula in question is known to hold modulo torsion. For smooth quasi-projective schemes over a field @@M16@@, the De Rham or Hodge version of formula (3.1) would moreover follow from the cycle-class version of formula (3.1) mod. algebraic equivalence, so that the De Rham or Hodge cohomological version gives another accessible way to test the validity of the ``cycle-class'' version of (3.1).

\sgaStableTarget{sga6.e14.para.3-5}\textbf{3.5.} The proof [1] of the transcendental cohomological version of (3.1) is essentially transcendental in nature, using tubular neighbourhoods of @@M17@@ in @@M18@@. This suggests that a proof of a formula of type (3.1) might be connected with the systematic consideration of groups @@M19@@ and cohomology groups \uline{with support in} @@M20@@, which might provide a sufficient technical means of ``localization around @@M21@@''.

\sgaStableTarget{sga6.e14.sec.4}\section*{4. Relations between @@M22@@ and the Chow ring @@M23@@}

\sgaStableTarget{sga6.e14.para.4-1}\textbf{4.1.} For simplicity, suppose that @@M24@@ is a smooth quasi-projective scheme over a field @@M25@@ (cf no.~8 for a more general case), so that we
